\documentclass{article}
\usepackage{graphicx}
\input{preamble}

\title{Assignment 4. \\ 3D Gaussian Splatting}
\date{}

\begin{document}

\maketitle

\section*{Assignment Description}

The goal of this assignment is to get familiar with the modern 3D scene reconstruction
method called \textbf{3D Gaussian Splatting} (3DGS).
Students will go through the full workflow from photographs to an interactive 3D viewer,
first on a reference dataset and then on a self-captured scene.
\\[4pt]
The assignment is worth \textbf{25~pts.} in total and consists of the following tasks:

\begin{enumerate}

% ── TASK 1 ───────────────────────────────────────────────────────────────────
\item \textbf{[5~pts.]} Running the workflow on a reference dataset.\\
This task is designed to help you \textbf{understand the full workflow} and verify
that your environment works correctly before you start working with your own dataset.
Working with reference data lets you learn each step without having to solve
capture-related problems at the same time.\\[4pt]
Download one of the publicly available datasets (e.g.\ \textit{Tanks and
Temples -- truck}, \textit{garden}, or another compatible dataset).
Use \textbf{only the photographs} from the dataset (the \texttt{input/} directory) --
\emph{no pre-computed data}.
Run the full workflow from scratch following the guide in \texttt{GUIDE.md}:

\begin{enumerate}
    \item \texttt{./run.sh} -- start the Docker container.
    \item \texttt{colmap gui} -- launch COLMAP GUI and perform:
        \begin{itemize}
            \item Feature extraction (camera model: \texttt{SIMPLE\_PINHOLE})
            \item Feature matching (Exhaustive)
            \item Reconstruction (Sparse)
            \item Export model to \texttt{distorted/sparse/0/}
        \end{itemize}
    \item \texttt{python3 convert.py -s data/<dataset>} -- image undistortion.
    \item \texttt{python3 train.py -s data/<dataset>} -- train the 3DGS model
          (30\,000 iterations, checkpoints at 7\,000 and 30\,000).
    \item \texttt{SIBR\_gaussianViewer\_app -m output/<uuid>} -- visualise the result.
\end{enumerate}

Document each step with screenshots (COLMAP GUI, training progress, viewer).
Report the \textbf{PSNR} metric at 7\,000 and 30\,000 iterations.

% ── TASK 2 ───────────────────────────────────────────────────────────────────
\item \textbf{[10~pts.]} Capturing your own scene with the Ximea camera.\\
Choose and capture your own scene using the industrial Ximea camera (high resolution).
\textbf{Downscale} the images before running COLMAP
(recommended: longer side to max.\ \textbf{600--1200~px}) -- the camera produces
images with very high resolution that would make both COLMAP and 3DGS training
impractically slow.

Run the full workflow in the same way as in Task~1.
Report the achieved PSNR and visually compare the result with the reference dataset.

% ── TASK 3 ───────────────────────────────────────────────────────────────────
\item \textbf{[5~pts.]} Evaluation of results.\\
Compare the results of both scenes (reference dataset vs.\ your own scene):
\begin{itemize}
    \item PSNR table at 7\,000 and 30\,000 iterations,
    \item visual comparison of renders from the same viewpoints,
    \item commentary -- what affected the quality of your scene
          (number of photos, lighting, object type, \ldots).
\end{itemize}

% ── TASK 4 ───────────────────────────────────────────────────────────────────
\item \textbf{[5~pts.]} Documentation.\\
Submit a report (PDF) containing:
\begin{itemize}
    \item description of each step with corresponding screenshots,
    \item description of the scene capture process (what, where, how),
    \item visual comparison of results,
    \item conclusion with recommendations -- what you would do differently for better results.
\end{itemize}

\end{enumerate}

\textbf{Submission deadline:} 12th lab session of the semester

\newpage

% =============================================================================
\section*{Notes, Tips and Resources}

\subsection*{Follow GUIDE.md}

All specific commands, paths and parameters are described in detail in
\texttt{GUIDE.md} in the repository.
If you encounter any problem, first read the \emph{Troubleshooting} section of the guide.

% ── Dataset ──────────────────────────────────────────────────────────────────
\subsection*{Where to find a reference dataset}

\begin{itemize}
    \item \textbf{Tanks and Temples} (truck, train):
          \href{https://www.tanksandtemples.org/download/}{tanksandtemples.org/download}
    \item \textbf{Mip-NeRF 360} (garden, room, counter, \ldots):
          \href{https://jonbarron.info/mipnerf360/}{jonbarron.info/mipnerf360}
    \item Use \textbf{only the folder with photographs} from the dataset --
          remove any pre-computed sparse/dense data.
\end{itemize}

% ── Downscale ─────────────────────────────────────────────────────────────────
\subsection*{Why and how to downscale images}

The Ximea camera produces images with resolutions up to 12--20~Mpx.
COLMAP extracts features at pixel level and at high resolution:
\begin{itemize}
    \item feature extraction can take many hours,
    \item GPU memory usage during training is significantly higher.
\end{itemize}
\textbf{Recommended resolution:} longer side at \textbf{1600--2000~px}
(approximately 2--3~Mpx).
Preserve the aspect ratio, do not crop.
Run the downscale \emph{before} starting COLMAP.

% ── Capture tips ─────────────────────────────────────────────────────────────
\subsection*{Tips for capturing a high-quality scene}

\begin{itemize}
    \item \textbf{Number of photographs:} at least \textbf{60--120 images}.
          More photos = better reconstruction, but longer COLMAP and training time.
    \item \textbf{Coverage:} move \emph{around the entire object} in a circle,
          ideally at several heights (low, middle, slightly from above).
          Every point in the scene must be visible from \textbf{at least 3 different angles}.
    \item \textbf{Image overlap:} rotate the camera by at most 10--15° between shots --
          adjacent images must share at least 60\,\% of the scene.
    \item \textbf{Static scene:} nothing must move during the entire capture
          (no people in frame, no wind, remove any moveable objects).
    \item \textbf{Lighting:} use uniform, diffuse light (overcast day or
          indoor scene with multiple light sources).
          Avoid harsh shadows and direct sunlight -- shadows change as the
          camera moves and confuse the reconstruction.
    \item \textbf{Object type:} the ideal object has a \emph{textured, matte}
          surface (e.g.\ a statue, furniture, basket, plant).
          \textbf{Problematic} surfaces:
          \begin{itemize}
              \item shiny / mirror-like surfaces (metal, glass),
              \item homogeneous surfaces without texture (white wall, plain table),
              \item transparent objects (glass, plastic).
          \end{itemize}
    \item \textbf{Camera settings:} use fixed settings (exposure, ISO, focus)
          throughout the entire capture session.
          Auto-focus changes between shots cause inconsistency.
    \item \textbf{Motion blur:} all photographs must be sharp.
          In low-light conditions, increase ISO rather than extending exposure time.
\end{itemize}

% ── Scene selection ──────────────────────────────────────────────────────────
\subsection*{Recommended scene types}

\begin{itemize}
    \item Tabletop object (figurine, plant, mug) on a neutral background --
          easy to walk around, short capture time.
    \item Room corner with furniture -- larger volume, more interesting result.
    \item Staircase or corridor -- good texture, but requires more images.
\end{itemize}

% ── Metrics ──────────────────────────────────────────────────────────────────
\subsection*{Interpreting the PSNR metric}

\begin{center}
\begin{tabular}{ll}
\toprule
PSNR & Quality \\
\midrule
$< 22$~dB & poor reconstruction (too few images, bad coverage) \\
$22$--$26$~dB & acceptable (checkpoint at 7\,000 it.) \\
$26$--$29$~dB & good (checkpoint at 30\,000 it.) \\
$> 29$~dB & excellent \\
\bottomrule
\end{tabular}
\end{center}

% ── Useful links ─────────────────────────────────────────────────────────────
\subsection*{Useful links}

\begin{itemize}
    \item Original paper: Kerbl et al., \textit{3D Gaussian Splatting for
          Real-Time Radiance Field Rendering}, SIGGRAPH 2023.\\
          \href{https://repo-sam.inria.fr/fungraph/3d-gaussian-splatting/}{repo-sam.inria.fr/fungraph/3d-gaussian-splatting}
    \item Visual tutorial:
          \href{https://learnopencv.com/3d-gaussian-splatting/}{learnopencv.com/3d-gaussian-splatting}
    \item COLMAP documentation:
          \href{https://colmap.github.io/}{colmap.github.io}
    \item \textbf{Repository with Docker environment, guides and assignment:}\\
          \href{https://github.com/michKovac/PVSO-gaussian-splatting-assignment}{github.com/michKovac/PVSO-gaussian-splatting-assignment}\\
          Contains: \texttt{GUIDE.md}, \texttt{Dockerfile}, \texttt{run.sh}, \texttt{run.bat}
\end{itemize}

\end{document}
